% UNIFIED 2026-05-13
% SOURCE MAP:
%   - Plan 1/faber-krahn-transfer.tex          (Kohler--Jobin transfer; constants c_1, c_2)
%   - Plan 1/faber-krahn-transfer.md           (same, prose version)
%   - Final/BoundedReduction.tex               (BDV truncation; c_SV(n,R) -> c_SV^{glob}(n))
%   - Final/KohlerJobinTransfer.tex            (canonical scale-invariant statement)
%   - Manuscript/00_notation_register.md       (n, R_*(n), delta_T, delta_FK, A=Fraenkel)
%   - Manuscript/source_map.md  §D.4           (this subsection's items D4.1-D4.6)
%
% Inputs consumed from earlier subsections of 02_plan1_draft.tex:
%   - Agent 7 (NearlySphericalClosure / SaintVenantAssembly), Theorem (SV-bounded):
%        E(\Om)-E(B_1) \ge c_{\rm SV}(n,R)\,\Fr(\Om)^2
%     for every open \Om \subset B_R with |\Om|=\omega_n, where R\ge 2.
%
% Notation register:
%   - dimension is n (N converted);
%   - \Fr := Fraenkel asymmetry (\mathcal A);
%   - \deltaT(\Om) := E(\Om)-E(B^*) (BDV-normalised, E=-T/2);
%   - \deltaFK(\Om) := (|\Om|/\omega_n)^{2/n}(\lambda_1(\Om)-\lambda_1(B^*));
%   - L_n := \lambda_1(B_1), T_n := T(B_1) = \omega_n/(n(n+2)).
%
% Output of this subsection:
%   - Theorem (BR): global sharp Saint--Venant, c_{\rm SV}^{\rm glob}(n);
%   - Theorem (KJ-FK): \Fr(\Om)^2 \le C_{\rm FK}(n)\,\deltaFK(\Om).
%
% This file is intended to be the final subsection of Manuscript/02_plan1_draft.tex
% and is integrated by Agent 19.

\subsection{Bounded reduction and the Kohler--Jobin transfer to Faber--Krahn}
\label{sec:plan1-kj-transfer}

This subsection closes Plan~1. The input is the bounded sharp
Saint--Venant inequality of Agent~7 (Theorem~\ref{thm:SV-bounded}).
Three steps follow:
\begin{enumerate}
\item[(i)] \emph{Bounded reduction}
(\S\ref{ssec:bounded-reduction}). Using the constructive BDV
truncation lemma, we replace an arbitrary open
$\Om\subset\R^n$ with finite measure by a translate that lies in a
ball of \emph{dimensional} radius $R_*(n)$, while controlling the loss
of both the Saint--Venant deficit and the Fraenkel asymmetry. This
removes the $R$-dependence from the Saint--Venant constant.
\item[(ii)] \emph{Kohler--Jobin inequality}
(\S\ref{ssec:kj-inequality}). We state the scale-invariant comparison
$T(\Om)\,\lambda_1(\Om)^{(n+2)/2}\ge T_n L_n^{(n+2)/2}$ and derive the
pointwise lower bound for $\lambda_1(\Om)-L$ in terms of the torsion
deficit.
\item[(iii)] \emph{Final Plan~1 theorem}
(\S\ref{ssec:plan1-final}). Combining (i) and (ii) yields the sharp
quantitative Faber--Krahn inequality in the canonical normalisation of
the notation register,
\[
\Fr(\Om)^2 \le C_{\rm FK}(n)\,\deltaFK(\Om),
\]
with $C_{\rm FK}(n)$ depending only on the dimension.
\end{enumerate}

\subsubsection*{Standing input from Plan~1, Agent~7.}

\begin{theorem}[Bounded sharp Saint--Venant; Agent~7, item D3.8]
\label{thm:SV-bounded}
For every $n\ge 2$ and every $R\ge 2$ there is a constant
$c_{\rm SV}(n,R)>0$, explicit in the BDV constants of the
selection / graph entry / nearly-spherical chain, such that
\begin{equation}\label{eq:SV-bounded}
\deltaT(\Om) \;=\; E(\Om)-E(B_1) \;\ge\; c_{\rm SV}(n,R)\,\Fr(\Om)^2
\end{equation}
for every open set $\Om\subset B_R$ with $|\Om|=\omega_n$.
\end{theorem}

The hypothesis $|\Om|=\omega_n$ is the standing volume normalisation
of the chain; we will remove it by scaling at the end
(Remark~\ref{rem:scale-removal}).

% =====================================================================
\subsubsection{Bounded reduction: removing the $R$-dependence}
\label{ssec:bounded-reduction}
% =====================================================================

The bounded reduction is the BDV truncation argument
\cite[Lemmas~5.1--5.3]{BDV}. It is constructive: a De~Giorgi
$L^\infty$ iteration is combined with an energy cut-off competitor;
no compactness, no minimising sequence, and no contradiction step is
used. The end product is a bounded surrogate $\Otilde$ of $\Om$ inside
a dimensional ball, with quantitatively controlled loss.

\paragraph{Scale-invariant Saint--Venant deficit.} It is convenient to
work with the scale-invariant deficit
\begin{equation}\label{eq:Dscaled}
D(\Om) \;:=\; \omega_n^{-(n+2)/n}\,\deltaT(\Om)
\;=\; \omega_n^{-(n+2)/n}\,\bigl(E(\Om)-E(B^*)\bigr).
\end{equation}
When $|\Om|=\omega_n$ we have $B^*=B_1$ and $D(\Om)\ge 0$ is the
volume-normalised Saint--Venant deficit. The Fraenkel asymmetry
$\Fr(\Om)$ is dimensionless and scale invariant.

\paragraph{Two BDV black boxes.} We import two estimates from
\cite{BDV} as black boxes.

\begin{lemma}[\protect{\cite[Lemma~5.1]{BDV}}: De~Giorgi $L^\infty$ tail]
\label{lem:51}
There is $C_7=C_7(n)>0$ such that, for every open $\Om$ with
$|\Om|=\omega_n$ and every $R\ge 1$,
\begin{equation}\label{eq:tail}
\|u_\Om\|_{L^\infty(\Om\setminus B_{R+1})}
\;\le\; C_7\,|\Om\setminus B_R|^{1/n}.
\end{equation}
\end{lemma}

\begin{proof}[Reference]
\cite[Lemma~5.1]{BDV}; De~Giorgi iteration on the truncated test
function $\varphi_k^2(u_\Om-s_k)_+$. Constructive; constant $C_7(n)$
explicit.
\end{proof}

\begin{lemma}[\protect{\cite[Lemma~5.2]{BDV}}: suboptimal
$\Fr^4$ bound]\label{lem:52}
There is $C_8=C_8(n)>0$ such that for every open $\Om$ of finite
measure,
\begin{equation}\label{eq:far}
D(\Om) \;\ge\; \frac{1}{C_8}\,\Fr(\Om)^4.
\end{equation}
\end{lemma}

\begin{proof}[Reference]
\cite[Lemma~5.2]{BDV}; equivalent, via $E=-T/2$, to
\cite[Theorem~1]{FMPLaplace} with $p=2,q=1$. Pure rearrangement plus
P\'olya--Szeg\H{o}; no contradiction.
\end{proof}

\paragraph{The truncation lemma.} The core construction.

\begin{lemma}[\protect{\cite[Lemma~5.3]{BDV}}: bounded truncation]
\label{lem:53}
There exist dimensional constants $C_9=C_9(n)>0$,
$\delta_{\rm BDV}=\delta_{\rm BDV}(n)\in(0,1]$, and
$d(n)\ge 1$ such that for every open $\Om$ with $|\Om|=\omega_n$ and
$D(\Om)\le\delta_{\rm BDV}(n)$, there exists an open set
$\Otilde\subset\R^n$ satisfying
\begin{align}
|\Otilde| &= \omega_n, \label{eq:53vol} \\
\mathrm{diam}(\Otilde) &\le d(n), \label{eq:53diam} \\
D(\Otilde) &\le C_9(n)\,D(\Om), \label{eq:53D} \\
\Fr(\Om) &\le \Fr(\Otilde) + C_9(n)\,D(\Om). \label{eq:53A}
\end{align}
\end{lemma}

\begin{proof}[Sketch and reference]
\cite[Lemma~5.3]{BDV}. The proof combines~\eqref{eq:tail} with the
cut-off test function $\varphi_k u_\Om$ in the Saint--Venant
variational characterisation of $E(\Om\cap B_{k+2})$, yielding
\cite[(5.11)]{BDV}
\[
E(\Om\cap B_{k+2}) \;\le\; E(\Om) + C\,b_k^{1+1/n},
\qquad b_k := |\Om\setminus B_k|/\omega_n.
\]
A geometric iteration on $k\le K\le K(n)$ produces $b_{K+1}\le
2\widehat C\,D(\Om)$. Set $r:=(|\Om\cap B_{K+3}|/\omega_n)^{1/n}\in
[1-\widehat C\,D(\Om),1]$ and $\Otilde:=r^{-1}(\Om\cap B_{K+3})$.
Properties~\eqref{eq:53vol}--\eqref{eq:53D} follow directly. The
asymmetry estimate~\eqref{eq:53A} uses
\[
|\Om\,\Delta\,(B_1+x_0)|
\;\le\;
|(\Om\cap B_{K+3})\,\Delta\,\Om|
+ r^n\,|\Otilde\,\Delta\,(B_1+x_0/r)|
+ \omega_n(1-r^n),
\]
followed by $1-r^n\le n\widehat C\,D(\Om)$. The construction is
\emph{constructive}: no minimising sequence, no Hausdorff limit, no
contradiction step.
\end{proof}

\begin{remark}[Translation freedom]\label{rem:translation}
The diameter bound~\eqref{eq:53diam} is translation invariant. After
translating $\Otilde$ we may assume $\Otilde\subset B_{d(n)}$. Both
$\Fr$ and $D$ are translation invariant.
\end{remark}

\paragraph{Working radius.} Define the \emph{bounded-reduction radius}
\begin{equation}\label{eq:Rstar}
R_*(n) \;:=\; \max\bigl\{d(n),\,2\bigr\}.
\end{equation}
By Remark~\ref{rem:translation} any $\Otilde$ produced by
Lemma~\ref{lem:53} can be translated into $B_{R_*(n)}$, so
Theorem~\ref{thm:SV-bounded} applies with the dimensional choice
$R=R_*(n)$.

\paragraph{Far-deficit branch.}

\begin{lemma}[Far-deficit quadratic bound]\label{lem:far}
For every open $\Om$ with $|\Om|=\omega_n$ and
$D(\Om)\ge\delta_{\rm BDV}(n)/4$,
\begin{equation}\label{eq:far-quad}
E(\Om)-E(B_1) \;\ge\;
\omega_n^{(n+2)/n}\,\frac{\delta_{\rm BDV}(n)}{16}\,\Fr(\Om)^2.
\end{equation}
\end{lemma}

\begin{proof}
Since $\Fr(\Om)<2$, $\Fr(\Om)^2/4<1$. Hence
$D(\Om)\ge\delta_{\rm BDV}(n)/4\ge(\delta_{\rm BDV}(n)/16)\Fr(\Om)^2$.
Multiplying by $\omega_n^{(n+2)/n}$ and using
$E(\Om)-E(B_1)=\omega_n^{(n+2)/n}D(\Om)$ at volume $\omega_n$ gives the
claim.
\end{proof}

\begin{remark}
The estimate~\eqref{eq:far-quad} is much weaker than~\eqref{eq:far}
(it converts the $\Fr^4$ control into an $\Fr^2$ control by absorbing
$\Fr^2\le 4$), but in this regime $D$ is bounded below by a
dimensional constant, so the trade is harmless.
\end{remark}

\paragraph{Near-deficit branch.}

\begin{lemma}[Near-deficit transfer]\label{lem:near}
Set
\begin{equation}\label{eq:csv-tilde}
\widetilde c_{\rm SV}(n)
\;:=\; \omega_n^{-(n+2)/n}\,c_{\rm SV}\bigl(n,R_*(n)\bigr)\;>\;0.
\end{equation}
For every open $\Om$ with $|\Om|=\omega_n$ and
$D(\Om)\le\delta_{\rm BDV}(n)$,
\begin{equation}\label{eq:near-Fr}
\Fr(\Om)^2 \;\le\;
\Bigl(\frac{2C_9(n)}{\widetilde c_{\rm SV}(n)}
+2C_9(n)^2\,\delta_{\rm BDV}(n)\Bigr)\,D(\Om).
\end{equation}
\end{lemma}

\begin{proof}
Apply Lemma~\ref{lem:53} to obtain $\Otilde$ satisfying
\eqref{eq:53vol}--\eqref{eq:53A}, and by Remark~\ref{rem:translation}
assume $\Otilde\subset B_{d(n)}\subset B_{R_*(n)}$. Apply
Theorem~\ref{thm:SV-bounded} with $R=R_*(n)$ to $\Otilde$:
\[
E(\Otilde)-E(B_1) \;\ge\; c_{\rm SV}\bigl(n,R_*(n)\bigr)\,\Fr(\Otilde)^2.
\]
Divide by $\omega_n^{(n+2)/n}$ (and use $|\Otilde|=\omega_n$) to obtain
\begin{equation}\label{eq:Dtilde-Atilde}
D(\Otilde) \;\ge\; \widetilde c_{\rm SV}(n)\,\Fr(\Otilde)^2.
\end{equation}
Combined with~\eqref{eq:53D},
\[
\Fr(\Otilde)^2 \;\le\;
\frac{D(\Otilde)}{\widetilde c_{\rm SV}(n)}
\;\le\;
\frac{C_9(n)\,D(\Om)}{\widetilde c_{\rm SV}(n)}.
\]
Hence
$\Fr(\Otilde)\le\bigl(C_9(n)D(\Om)/\widetilde c_{\rm SV}(n)\bigr)^{1/2}$.
Plugging this into~\eqref{eq:53A} and squaring,
\begin{align*}
\Fr(\Om)^2
&\le \bigl(\Fr(\Otilde) + C_9(n)\,D(\Om)\bigr)^2
\;\le\; 2\Fr(\Otilde)^2 + 2\,C_9(n)^2\,D(\Om)^2 \\
&\le \frac{2C_9(n)}{\widetilde c_{\rm SV}(n)}\,D(\Om)
   + 2\,C_9(n)^2\,D(\Om)^2 \\
&\le \Bigl(\frac{2C_9(n)}{\widetilde c_{\rm SV}(n)}
   + 2\,C_9(n)^2\,\delta_{\rm BDV}(n)\Bigr)\,D(\Om),
\end{align*}
where the last step uses $D(\Om)\le\delta_{\rm BDV}(n)\le 1$.
\end{proof}

\begin{corollary}[Near-deficit quadratic bound]\label{cor:near}
With
\begin{equation}\label{eq:c-near-glob}
c_{\rm near}^{\rm glob}(n)
\;:=\;
\frac{\omega_n^{(n+2)/n}}
{2C_9(n)\bigl(1/\widetilde c_{\rm SV}(n)+C_9(n)\delta_{\rm BDV}(n)\bigr)},
\end{equation}
every open $\Om$ with $|\Om|=\omega_n$ and
$D(\Om)\le\delta_{\rm BDV}(n)$ satisfies
\begin{equation}\label{eq:near-final}
E(\Om)-E(B_1) \;\ge\; c_{\rm near}^{\rm glob}(n)\,\Fr(\Om)^2.
\end{equation}
\end{corollary}

\begin{proof}
\eqref{eq:near-Fr} gives
$D(\Om)\ge\Fr(\Om)^2/\bigl(2C_9/\widetilde c_{\rm SV}+2C_9^2\delta_{\rm BDV}\bigr)$.
Multiply by $\omega_n^{(n+2)/n}$ and use $E-E(B_1)=\omega_n^{(n+2)/n}D$
to get~\eqref{eq:near-final}.
\end{proof}

\paragraph{Global bounded reduction.}

\begin{theorem}[Global sharp Saint--Venant; bounded reduction]
\label{thm:BR}
Define
\begin{equation}\label{eq:csv-glob}
c_{\rm SV}^{\rm glob}(n)
\;:=\;
\min\!\Bigl\{
c_{\rm near}^{\rm glob}(n),\;
\omega_n^{(n+2)/n}\,\delta_{\rm BDV}(n)/16
\Bigr\}\;>\;0.
\end{equation}
Then for every open set $\Om\subset\R^n$ with $|\Om|=\omega_n$,
\begin{equation}\label{eq:BR}
\deltaT(\Om) \;=\; E(\Om)-E(B_1)
\;\ge\; c_{\rm SV}^{\rm glob}(n)\,\Fr(\Om)^2.
\end{equation}
The constant $c_{\rm SV}^{\rm glob}(n)$ depends only on the dimension
$n$: the $R$-dependence inherited from
Theorem~\ref{thm:SV-bounded} has been internalised at the
dimensional choice $R=R_*(n)$.
\end{theorem}

\begin{proof}
The bound $D(\Om)\ge 0$ is the classical Saint--Venant inequality.
We split into two cases.

\emph{Case A} ($D(\Om)\ge\delta_{\rm BDV}(n)/4$):
Lemma~\ref{lem:far} gives
$E(\Om)-E(B_1)\ge\omega_n^{(n+2)/n}\bigl(\delta_{\rm BDV}(n)/16\bigr)\Fr(\Om)^2$,
which is~\eqref{eq:BR} with constant
$\omega_n^{(n+2)/n}\delta_{\rm BDV}(n)/16$.

\emph{Case B} ($D(\Om)\le\delta_{\rm BDV}(n)$, which in particular
covers $D(\Om)\le\delta_{\rm BDV}(n)/4$):
Corollary~\ref{cor:near} applies and gives~\eqref{eq:BR} with
constant $c_{\rm near}^{\rm glob}(n)$.

The two cases overlap on
$\delta_{\rm BDV}(n)/4\le D(\Om)\le\delta_{\rm BDV}(n)$ and together
cover $D(\Om)\in[0,+\infty)$. Taking the minimum of the two
constants gives~\eqref{eq:csv-glob}.
\end{proof}

\begin{remark}[Dependence on $R$]\label{rem:R-dropout}
The dependence on the confining radius $R$ enters only via
Theorem~\ref{thm:SV-bounded}, which is applied to the surrogate
$\Otilde\subset B_{R_*(n)}$. Since $R_*(n)$ is itself dimensional, the
final constant $c_{\rm SV}^{\rm glob}(n)$ in~\eqref{eq:csv-glob}
depends only on $n$. Concretely, the
$R$-dependence is absorbed in
$\widetilde c_{\rm SV}(n)=\omega_n^{-(n+2)/n}c_{\rm SV}(n,R_*(n))$ in
\eqref{eq:csv-tilde}; everything downstream is purely dimensional.
\end{remark}

\begin{remark}[Removal of the volume normalisation]\label{rem:scale-removal}
\eqref{eq:BR} stated for $|\Om|=\omega_n$ extends to every open
$\Om\subset\R^n$ with $0<|\Om|<\infty$ by scale invariance of $D$ and
$\Fr$: writing
$\Om_*:=(\omega_n/|\Om|)^{1/n}\Om$ one has $|\Om_*|=\omega_n$,
$\Fr(\Om_*)=\Fr(\Om)$, and $D(\Om_*)=D(\Om)$. Hence~\eqref{eq:BR}
applied to $\Om_*$ reads
\begin{equation}\label{eq:BR-scaleinv}
\omega_n^{(n+2)/n}\,D(\Om) \;=\; |\Om|^{(n+2)/n}\,D(\Om)\cdot
\bigl(\omega_n/|\Om|\bigr)^{(n+2)/n}
\quad\text{etc.}
\end{equation}
Equivalently, at volume $|B^*|=|\Om|$,
\begin{equation}\label{eq:torsion-deficit-out}
T(B^*) - T(\Om)
\;=\; 2\bigl(E(\Om)-E(B^*)\bigr)
\;\ge\; 2\,c_{\rm SV}^{\rm glob}(n)\,
\Bigl(\frac{|\Om|}{\omega_n}\Bigr)^{(n+2)/n}\Fr(\Om)^2.
\end{equation}
This is the input to the Kohler--Jobin transfer
(\S\ref{ssec:kj-inequality}).
\end{remark}

% =====================================================================
\subsubsection{The Kohler--Jobin inequality and the pointwise transfer}
\label{ssec:kj-inequality}
% =====================================================================

\paragraph{Statement.} We invoke the Kohler--Jobin inequality in the
following scale-invariant form, which is the form needed to transfer
Saint--Venant stability to Faber--Krahn stability.

\begin{theorem}[Kohler--Jobin inequality]\label{thm:kj}
For every open set $\Om\subset\R^n$ with $0<|\Om|<\infty$ and
$\lambda_1(\Om)<\infty$, and for every ball $B$,
\begin{equation}\label{eq:kj-global}
T(\Om)\,\lambda_1(\Om)^{(n+2)/2}
\;\ge\; T(B)\,\lambda_1(B)^{(n+2)/2}.
\end{equation}
Equivalently, $\Psi(\Om):=T(\Om)\,\lambda_1(\Om)^{(n+2)/2}\ge
\Psi(B_1)=T_n L_n^{(n+2)/2}$. Equality holds if and only if $\Om$ is a
ball (up to translation and a Lebesgue null set).
\end{theorem}

\begin{proof}[Reference]
The original statement is in Kohler--Jobin \cite{KohlerJobin1978}. We
use the form of \cite[Theorem~3.3, p.~322]{Brasco2014}, proved by
Schwarz symmetrisation and a one-dimensional rearrangement of the
torsion potential along level sets of the eigenfunction. The
characterisation of the equality case is \cite[Theorem~3.3]{Brasco2014}.
\end{proof}

\begin{remark}[Canonical normalisation and scaling]\label{rem:kj-canonical}
The exponent $(n+2)/2$ in~\eqref{eq:kj-global} is forced by scale
invariance: from $T(t\Om)=t^{n+2}T(\Om)$ and
$\lambda_1(t\Om)=t^{-2}\lambda_1(\Om)$ we read
$\Psi(t\Om)=t^{n+2-2\cdot(n+2)/2}\Psi(\Om)=\Psi(\Om)$, so $\Psi$ is
scale-invariant. This is the canonical form of the inequality used
throughout the manuscript; both source files
\texttt{Plan~1/faber-krahn-transfer.tex} and
\texttt{Final/KohlerJobinTransfer.tex} use the same form
(with $N$ converted to $n$).
\end{remark}

\paragraph{Pointwise transfer at fixed volume.} Fix the comparison ball
$B^*$ with $|B^*|=|\Om|$, and write
\[
L:=\lambda_1(B^*),\qquad T_0:=T(B^*),\qquad \beta:=\frac{n+2}{2}.
\]
By scaling $L=L_n(\omega_n/|B^*|)^{2/n}$ and
$T_0=T_n(|B^*|/\omega_n)^{(n+2)/n}$.

\begin{lemma}[Pointwise Kohler--Jobin transfer]\label{lem:KJ-ptwise}
For every open $\Om$ with $|\Om|=|B^*|$ and $\lambda_1(\Om)<\infty$,
\begin{equation}\label{eq:plan1-KJ-ptwise}
\lambda_1(\Om)
\;\ge\; L\cdot\Bigl(\frac{T_0}{T(\Om)}\Bigr)^{1/\beta}
\;=\; L\cdot\Bigl(\frac{T_0}{T(\Om)}\Bigr)^{2/(n+2)}.
\end{equation}
\end{lemma}

\begin{proof}
Apply Theorem~\ref{thm:kj} with $B=B^*$ and rearrange. Since $T(\Om)>0$
and $\lambda_1(\Om)<\infty$, both sides are finite and positive.
\end{proof}

\begin{lemma}[Elementary one-variable inequality]\label{lem:plan1-onevar}
For every $p\in(0,1]$ and every $s\in[0,1)$,
\begin{equation}\label{eq:plan1-onevar}
(1-s)^{-p}-1 \;\ge\; p\,s.
\end{equation}
\end{lemma}

\begin{proof}
Set $f(s):=(1-s)^{-p}-1$. Then $f(0)=0$ and
$f'(s)=p(1-s)^{-(p+1)}\ge p$ for $s\in[0,1)$. By the fundamental
theorem of calculus,
$f(s)=\int_0^s f'(\tau)\,d\tau\ge ps$, which is~\eqref{eq:plan1-onevar}.
The factor $p$ on the right is sharp: $f'(0)=p$, so the linear
approximation is tight at $s=0$.
\end{proof}

\paragraph{Small-deficit linearisation.} The classical Saint--Venant
inequality gives $T(\Om)\le T_0$ at fixed volume, so the torsion
deficit $\sigma(\Om):=T_0-T(\Om)\ge 0$. (We use the local symbol
$\sigma$ here, since the manuscript reserves $\delta_T$ for the
energy-form deficit.) Equivalently
$\sigma(\Om)=2\,\deltaT(\Om)$, and setting
$s:=\sigma(\Om)/T_0\in[0,1)$, \eqref{eq:plan1-KJ-ptwise} becomes
\begin{equation}\label{eq:KJ-s}
\lambda_1(\Om)-L \;\ge\; L\bigl[(1-s)^{-1/\beta}-1\bigr].
\end{equation}

\begin{lemma}[Small-deficit Kohler--Jobin linearisation]\label{lem:KJ-linear}
For every open $\Om$ with $|\Om|=|B^*|$ and $\sigma(\Om)\le T_0/2$,
\begin{equation}\label{eq:plan1-KJ-linear}
\lambda_1(\Om)-L
\;\ge\; \frac{2L}{(n+2)\,T_0}\,\sigma(\Om)
\;=\; \frac{4L}{(n+2)\,T_0}\,\deltaT(\Om).
\end{equation}
\end{lemma}

\begin{proof}
With $p=1/\beta=2/(n+2)\in(0,1]$ and $s=\sigma(\Om)/T_0\in[0,1/2]$,
combine~\eqref{eq:KJ-s} and Lemma~\ref{lem:plan1-onevar}:
\[
\lambda_1(\Om)-L
\;\ge\; L\,p\,s
\;=\; \frac{2L}{(n+2)\,T_0}\,\sigma(\Om).
\]
The second equality uses $\sigma=2\deltaT$.
\end{proof}

\begin{remark}[Universal multiplier]\label{rem:multiplier}
The factor $2L/((n+2)T_0)$ is the \emph{Kohler--Jobin multiplier}.
Both $L$ and $T_0$ scale, but the product is dimensionless:
$2L/((n+2)T_0)=2L_n/((n+2)T_n)\cdot(\omega_n/|B^*|)^{(n+2)/n}/(\omega_n/|B^*|)^{2/n}\cdots$,
which simplifies to
\begin{equation}\label{eq:multiplier-explicit}
\frac{2L}{(n+2)T_0}
\;=\;
\frac{2L_n}{(n+2)T_n}\cdot\Bigl(\frac{\omega_n}{|B^*|}\Bigr)^{(n+4)/n}\;.
\end{equation}
Using $T_n=\omega_n/(n(n+2))$, the dimensional part simplifies to
$2L_n/((n+2)T_n)=2nL_n/\omega_n$.
\end{remark}

\paragraph{Large-deficit branch.}

\begin{lemma}[Large-deficit Kohler--Jobin bound]\label{lem:KJ-large}
For every open $\Om$ with $|\Om|=|B^*|$ and $\sigma(\Om)>T_0/2$,
\begin{equation}\label{eq:KJ-large}
\lambda_1(\Om)-L
\;\ge\; L\bigl(2^{2/(n+2)}-1\bigr).
\end{equation}
\end{lemma}

\begin{proof}
$\sigma>T_0/2$ means $T_0/T(\Om)>2$, so~\eqref{eq:plan1-KJ-ptwise} gives
$\lambda_1(\Om)\ge L\cdot 2^{2/(n+2)}$, i.e.~\eqref{eq:KJ-large}.
\end{proof}

Since $|\Om|=|B^*|$ forces $\Fr(\Om)\in[0,2)$ (so $\Fr^2\le 4$),
in the large-deficit regime we may rewrite~\eqref{eq:KJ-large} as
\begin{equation}\label{eq:KJ-large-Fr}
\lambda_1(\Om)-L
\;\ge\; \frac{L\bigl(2^{2/(n+2)}-1\bigr)}{4}\,\Fr(\Om)^2.
\end{equation}

% =====================================================================
\subsubsection{Final Plan~1 sharp Faber--Krahn theorem}
\label{ssec:plan1-final}
% =====================================================================

We now combine the global Saint--Venant inequality
(Theorem~\ref{thm:BR}) and the pointwise Kohler--Jobin transfer
(Lemmas~\ref{lem:KJ-linear} and~\ref{lem:KJ-large}) to obtain the
sharp quantitative Faber--Krahn inequality in the canonical
normalisation of the notation register.

\paragraph{Statement in the lower form.}

\begin{theorem}[Sharp Faber--Krahn via Kohler--Jobin; lower form]
\label{thm:KJ-FK-lower}
There is a dimensional constant $c_{\rm FK}(n)>0$ such that, for
every open $\Om\subset\R^n$ with $0<|\Om|<\infty$,
\begin{equation}\label{eq:KJ-FK-lower}
\deltaFK(\Om)
\;=\;
\Bigl(\frac{|\Om|}{\omega_n}\Bigr)^{2/n}
\bigl(\lambda_1(\Om)-\lambda_1(B^*)\bigr)
\;\ge\; c_{\rm FK}(n)\,\Fr(\Om)^2.
\end{equation}
The constant is given explicitly by
\begin{equation}\label{eq:cFKdef}
c_{\rm FK}(n) \;:=\;
\min\!\Bigl\{
c_1(n),\; c_2(n)
\Bigr\},
\end{equation}
where
\begin{equation}\label{eq:c1c2-def}
c_1(n) := \frac{4 L_n\,c_{\rm SV}^{\rm glob}(n)}{(n+2)\,T_n}
\;=\; \frac{4 n L_n}{\omega_n}\,c_{\rm SV}^{\rm glob}(n),
\qquad
c_2(n) := \frac{L_n\bigl(2^{2/(n+2)}-1\bigr)}{4}.
\end{equation}
\end{theorem}

\begin{proof}
We first reduce to the volume-normalised case $|\Om|=\omega_n$. The
quantity $\deltaFK(\Om)$ is scale invariant: if
$\Om_*:=(\omega_n/|\Om|)^{1/n}\Om$, then $|\Om_*|=\omega_n$,
$\Fr(\Om_*)=\Fr(\Om)$, and
$\deltaFK(\Om_*)=\lambda_1(\Om_*)-L_n=
(|\Om|/\omega_n)^{2/n}(\lambda_1(\Om)-L)$ by the scaling
$\lambda_1(t\Om)=t^{-2}\lambda_1(\Om)$ (with $t=(\omega_n/|\Om|)^{1/n}$),
which equals $\deltaFK(\Om)$.
Thus it suffices to prove
$\lambda_1(\Om)-L_n\ge c_{\rm FK}(n)\Fr(\Om)^2$ for
$|\Om|=\omega_n$ (in which case $B^*=B_1$, $L=L_n$, $T_0=T_n$,
$\deltaFK(\Om)=\lambda_1(\Om)-L_n$).

Assume henceforth $|\Om|=\omega_n$. Theorem~\ref{thm:BR} gives
$\deltaT(\Om)\ge c_{\rm SV}^{\rm glob}(n)\Fr(\Om)^2$, so the torsion
deficit
\begin{equation}\label{eq:torsion-input}
\sigma(\Om) \;=\; T_n - T(\Om) \;=\; 2\deltaT(\Om)
\;\ge\; 2\,c_{\rm SV}^{\rm glob}(n)\,\Fr(\Om)^2.
\end{equation}

\emph{Case 1: $\sigma(\Om)\le T_n/2$.}
Lemma~\ref{lem:KJ-linear} (with $L=L_n$, $T_0=T_n$)
and~\eqref{eq:torsion-input} give
\[
\lambda_1(\Om) - L_n
\;\ge\;
\frac{2 L_n}{(n+2) T_n}\,\sigma(\Om)
\;\ge\;
\frac{2 L_n}{(n+2) T_n}\cdot 2\,c_{\rm SV}^{\rm glob}(n)\,\Fr(\Om)^2
\;=\;
\frac{4 L_n\,c_{\rm SV}^{\rm glob}(n)}{(n+2) T_n}\,\Fr(\Om)^2
\;=\;
c_1(n)\,\Fr(\Om)^2,
\]
the last equality being the definition of $c_1(n)$
in~\eqref{eq:c1c2-def}. The first inequality used the small-deficit
Kohler--Jobin linearisation expressed in terms of the torsion deficit
$\sigma$; the second inequality used $\sigma=2\deltaT$ combined with
the Saint--Venant input
$\deltaT\ge c_{\rm SV}^{\rm glob}\Fr^2$.

\emph{Case 2: $\sigma(\Om)>T_n/2$.}
\eqref{eq:KJ-large-Fr} with $L=L_n$ gives
\[
\lambda_1(\Om)-L_n
\;\ge\; \frac{L_n\bigl(2^{2/(n+2)}-1\bigr)}{4}\,\Fr(\Om)^2
\;=\; c_2(n)\,\Fr(\Om)^2.
\]

Combining the two cases gives
$\lambda_1(\Om)-L_n\ge\min\{c_1(n),c_2(n)\}\Fr(\Om)^2
=c_{\rm FK}(n)\Fr(\Om)^2$, which is~\eqref{eq:KJ-FK-lower} in the
volume-normalised case, and hence (by the first paragraph) in general.
\end{proof}

\begin{remark}[Canonical form of $c_1$]\label{rem:factor-two}
The constant $c_1(n)$ in~\eqref{eq:c1c2-def} matches
\cite[\S3, eq.~(3.1)]{Brasco2014} and the canonical form in
\texttt{Final/KohlerJobinTransfer.tex} (equations $c_1$, $c_{\rm FK}$).
Note that the source file
\texttt{Plan~1/faber-krahn-transfer.tex} states the small-deficit
multiplier as $2L/((n+2)T_0)$ applied to the torsion deficit
$\sigma=T_0-T$, while our proof applies it after the substitution
$\sigma=2\deltaT$; this accounts for the explicit factor of $4$ (not
$2$) on $L_n c_{\rm SV}^{\rm glob}/((n+2)T_n)$ in
\eqref{eq:c1c2-def}. The two source files
\texttt{Final/KohlerJobinTransfer.tex} (eq.~(\ref{eq:c1c2-def}) above)
and \texttt{Plan~1/faber-krahn-transfer.tex} (\S\,Quantitative
dependence) agree on this normalisation.
\end{remark}

\paragraph{Convenient abbreviation.} Throughout the rest of the
manuscript we write
\begin{equation}\label{eq:c1-canonical}
c_1(n) \;:=\; \frac{4 L_n\,c_{\rm SV}^{\rm glob}(n)}{(n+2)\,T_n}
\;=\; \frac{4 n L_n}{\omega_n}\,c_{\rm SV}^{\rm glob}(n),
\end{equation}
which is the definition just used.

\paragraph{Canonical statement (upper form).}
Inverting~\eqref{eq:KJ-FK-lower}, the inequality reads
$\Fr(\Om)^2\le C_{\rm FK}(n)\,\deltaFK(\Om)$ with
$C_{\rm FK}(n)=1/c_{\rm FK}(n)$. This is the upper form required by
the master brief (and shared with Plan~2).

\begin{theorem}[Sharp Faber--Krahn; Plan~1, upper form]
\label{thm:KJ-FK-upper}
There is a dimensional constant $C_{\rm FK}(n)>0$ such that, for every
open set $\Om\subset\R^n$ with $0<|\Om|<\infty$,
\begin{equation}\label{eq:KJ-FK-upper}
\boxed{\;
\Fr(\Om)^2 \;\le\; C_{\rm FK}(n)\,\deltaFK(\Om),
\;}
\end{equation}
where
\begin{equation}\label{eq:CFK}
C_{\rm FK}(n)
\;:=\; \frac{1}{c_{\rm FK}(n)}
\;=\; \max\!\Bigl\{
\frac{(n+2)\,T_n}{4 L_n\,c_{\rm SV}^{\rm glob}(n)},\;\;
\frac{4}{L_n\bigl(2^{2/(n+2)}-1\bigr)}
\Bigr\},
\end{equation}
with $c_{\rm SV}^{\rm glob}(n)$ defined in~\eqref{eq:csv-glob}.
Equality cases: by Theorem~\ref{thm:kj}, $\deltaFK(\Om)=0$
iff $\Om$ is a ball; in this case $\Fr(\Om)=0$.
\end{theorem}

\begin{proof}
If $\lambda_1(\Om)=+\infty$ then $\deltaFK(\Om)=+\infty$ and the
inequality \eqref{eq:KJ-FK-upper} is trivial (the left-hand side is
bounded above by $4$ since $\Fr(\Om)<2$). Otherwise
$\lambda_1(\Om)<\infty$ and the result is immediate from
Theorem~\ref{thm:KJ-FK-lower} with the canonical $c_1(n)$
from~\eqref{eq:c1-canonical}: invert the lower bound. The
constant $C_{\rm FK}(n)$ depends only on $n$ because
$c_{\rm SV}^{\rm glob}(n)$ does (Remark~\ref{rem:R-dropout}) and $L_n,
T_n$ are dimensional.
\end{proof}

\begin{remark}[$R$ drops out]\label{rem:R-removal-final}
The dependence on the confining radius $R$ is removed in two stages:
\begin{itemize}
\item Theorem~\ref{thm:SV-bounded} (Agent~7) provides
$c_{\rm SV}(n,R)$ for $\Om\subset B_R$. We \emph{instantiate} this
input at $R=R_*(n)=\max\{d(n),2\}$, a purely dimensional radius
defined in~\eqref{eq:Rstar}.
\item The bounded reduction (Theorem~\ref{thm:BR}) transfers this
to $c_{\rm SV}^{\rm glob}(n)$ via the BDV truncation lemma, which
produces a translate $\Otilde\subset B_{R_*(n)}$.
\end{itemize}
No subsequent step reintroduces $R$: the Kohler--Jobin inequality
\eqref{eq:kj-global} is independent of any confining ball, and the
pointwise transfer~\eqref{eq:plan1-KJ-ptwise} depends only on the
universal ball constants $L_n=\lambda_1(B_1)$, $T_n=T(B_1)$, and the
torsion deficit at fixed volume.
\end{remark}

\begin{remark}[Sharpness of the exponent]
The exponent $2$ on $\Fr$ in~\eqref{eq:KJ-FK-upper} is sharp; it is
saturated by smooth volume-preserving ellipsoidal perturbations of the
ball, for which $\Fr(\Om)$ and $\deltaFK(\Om)^{1/2}$ are of the same
order. See \cite[\S1.2]{BDV} for the construction.
\end{remark}

\begin{remark}[Match with Plan~2 normalisation]\label{rem:plan2-match}
The final theorem~\eqref{eq:KJ-FK-upper} is stated in the canonical
normalisation of the notation register (item~D4.6 of
\texttt{Manuscript/source\_map.md}), with $\deltaFK$ defined as in
notation register \S4: $\deltaFK(\Om)=(|\Om|/\omega_n)^{2/n}(\lambda_1(\Om)-\lambda_1(B^*))$.
This is exactly the normalisation of the Plan~2 final theorem
(item~E5.8); the two proofs share both the bounded reduction
(Theorem~\ref{thm:BR}) and the Kohler--Jobin transfer step
(\S\ref{ssec:kj-inequality}). The constant $C_{\rm FK}(n)$ produced
here and the constant $C_{\rm FK}^{\rm Plan~2}(n)$ produced in
\S\ref{sec:plan2-global-assembly} differ only through the
Saint--Venant constant $c_{\rm SV}^{\rm glob}(n)$, which each route
constructs independently.
\end{remark}

\paragraph{Summary of constants.}
\renewcommand{\arraystretch}{1.25}
\begin{center}
\begin{tabular}{|l|l|l|}
\hline
Symbol & Value/dependence & Source \\
\hline
$\omega_n$ & $|B_1|$ & dimensional \\
$T_n$ & $\omega_n/(n(n+2))$ & dimensional \\
$L_n$ & $\lambda_1(B_1)=j_{n/2-1,1}^2$ & dimensional \\
$\beta$ & $(n+2)/2$ & Theorem~\ref{thm:kj} \\
$C_7(n)$ & \cite[Lem.~5.1]{BDV} & Lemma~\ref{lem:51} \\
$C_8(n)$ & \cite[Lem.~5.2]{BDV} & Lemma~\ref{lem:52} \\
$C_9(n)$ & \cite[Lem.~5.3]{BDV} & Lemma~\ref{lem:53} \\
$\delta_{\rm BDV}(n)\in(0,1]$ & \cite[Lem.~5.3]{BDV} & Lemma~\ref{lem:53} \\
$d(n)\ge 1$ & \cite[Lem.~5.3]{BDV} & Lemma~\ref{lem:53} \\
$R_*(n)=\max\{d(n),2\}$ & dimensional & \eqref{eq:Rstar} \\
$c_{\rm SV}(n,R)$ & Agent~7 & Theorem~\ref{thm:SV-bounded} \\
$\widetilde c_{\rm SV}(n)$ & $\omega_n^{-(n+2)/n}c_{\rm SV}(n,R_*(n))$ & \eqref{eq:csv-tilde} \\
$c_{\rm near}^{\rm glob}(n)$ & \eqref{eq:c-near-glob} & Corollary~\ref{cor:near} \\
$c_{\rm SV}^{\rm glob}(n)$ & $\min\{c_{\rm near}^{\rm glob},\omega_n^{(n+2)/n}\delta_{\rm BDV}/16\}$ & Theorem~\ref{thm:BR} \\
$c_1(n)$ & $4L_n c_{\rm SV}^{\rm glob}/((n+2)T_n)=4nL_n c_{\rm SV}^{\rm glob}/\omega_n$ & \eqref{eq:c1-canonical} \\
$c_2(n)$ & $L_n(2^{2/(n+2)}-1)/4$ & \eqref{eq:c1c2-def} \\
$c_{\rm FK}(n)$ & $\min\{c_1(n),c_2(n)\}$ & Theorem~\ref{thm:KJ-FK-lower} \\
$C_{\rm FK}(n)$ & $1/c_{\rm FK}(n)$ & Theorem~\ref{thm:KJ-FK-upper} \\
\hline
\end{tabular}
\end{center}

\paragraph{Where compactness could have entered (and does not).}
The bounded reduction (\S\ref{ssec:bounded-reduction}) is the only step
where one might fear a hidden compactness argument. We record:
\begin{itemize}
\item BDV Lemma~5.1 (Lemma~\ref{lem:51}) is De~Giorgi
$L^\infty$ iteration: constructive, no minimising sequence.
\item BDV Lemma~5.2 (Lemma~\ref{lem:52}) is rearrangement plus
P\'olya--Szeg\H{o}: no contradiction.
\item BDV Lemma~5.3 (Lemma~\ref{lem:53}) reaches the truncation
index $K\le K(n)$ by a geometric iteration
\cite[(5.12)--(5.14)]{BDV}, not by Hausdorff compactness.
\item The gluing of near and far branches (Theorem~\ref{thm:BR}) is
arithmetic.
\item The Kohler--Jobin transfer (\S\ref{ssec:kj-inequality}) uses
only \eqref{eq:kj-global}, the FTC, and elementary monotone-function
inequalities.
\end{itemize}

This closes Plan~1.
